% 模型评价与推广
\section{模型评价与推广}
\label{sec:eval}

\subsection{模型优点}

\begin{enumerate}
\item \textbf{物理机理一致}：能耗采用题目统一的等效航程模型，即等效航程随
载荷的 1.5 次方收缩，飞行时间按爬升、巡航、下降三段解析计算，高程取自数字
高程模型；通信按自由空间损耗加逐点地形遮挡判定，全部参数可核验，结果可复现。
\item \textbf{双层校验闭环}：运输方案以硬约束加软约束双重检验，即首飞批与
医疗货箱时限为硬约束、普通批超期加权惩罚为软约束，零迟到方案才被采纳；
通信方案以 30\,m 步长逐点采样验证，保证飞行全程无盲区是确定性判定而非估算。
\item \textbf{协同设计实现双赢}：问题三把通信窗口作为运输重排的约束而非
惩罚，通过 10 项货箱重排同时减少架次、缩短完工、降低能耗并消除盲区，体现
约束也是优化资源的建模思想。
\item \textbf{资源核算诚实}：问题四对按组独立配置与集中池化做了定量对比，
明确指出库存约束下的最优组织方式，结论具有工程可操作性。
\end{enumerate}

\subsection{模型局限}

\begin{enumerate}
\item 未考虑随机因素，即阵风、突发禁飞与设备故障，调度为确定性模型；故障
冗余仅在问题四的静态备份层面讨论。
\item 中继为单点悬停，若中继无人机遇故障，其覆盖时段将出现盲区；规模更大
时可考虑双中继并机或移动中继。
\item 电池充电按两阶段曲线近似，未建模电池老化与温度影响，荷电状态均按
线性电量核算。
\item 模拟退火的随机性使结果存在微小波动，虽以固定随机种子保证复现，但无法
保证全局最优。
\end{enumerate}

\subsection{模型推广}

\begin{enumerate}
\item \textbf{动态重规划}：将模型嵌入滚动时域框架，随灾情变化，即新需求、
道路恢复与区域失联，在线重规划，即可支持多日持续保障。
\item \textbf{多中继自组织}：把布设点集扩展为可移动中继路径，模型可推广到
中继随运输进度迁移的动态覆盖问题。
\item \textbf{异构协同场景}：模型框架，即等效航程能耗、截止期优先调度与覆盖
采样三件套，可移植到海上搜救、森林消防、重大活动保障等大范围加弱通信的
应急场景。
\end{enumerate}